\chapter{泰勒斯定理}\label{ch:g9:thales}

不爬上去，你如何测量\omterm{def:g8:solids:def}{棱锥}的高度？米利都的泰勒斯
把它的影子与一根木棍的影子作了比较。
他的定理——\omterm{def:g6:lines:perp}{平行}线截得的\omterm{def:g6:lines:objects}{线段}成比例——
是比例模型、地图与放大图的数学核心，
也是最古老的有名字的定理之一。

\section{定理}

\begin{theorem}[泰勒斯]\label{thm:g9:thales:direct}
设两条直线交于点 $A$，在第一条直线上取 $M$，
在第二条直线上取 $N$，其中 $B$ 为第一条直线上较远的点，
$C$ 为第二条直线上较远的点。若直线 $(MN)$ 与 $(BC)$ \emph{\omterm{def:g6:lines:perp}{平行}}，则
三角形 $AMN$ 与 $ABC$ 的对应边成比例：
\[
\frac{AM}{AB} = \frac{AN}{AC} = \frac{MN}{BC}.
\]
\end{theorem}
\admitted

\begin{omfigure}
\begin{tikzpicture}[scale=1.05]
  \coordinate (A) at (0,0);
  \coordinate (B) at (5.0,0.6);
  \coordinate (C) at (3.2,2.8);
  \coordinate (M) at ($(A)!0.55!(B)$);
  \coordinate (N) at ($(A)!0.55!(C)$);
  \draw[thick] (A) -- (B) -- (C) -- cycle;
  \draw[omThm, very thick] (M) -- (N);
  \draw[omThm, very thick] (B) -- (C);
  \fill (A) circle (1.5pt) node[left, font=\small] {$A$};
  \fill (B) circle (1.5pt) node[right, font=\small] {$B$};
  \fill (C) circle (1.5pt) node[above, font=\small] {$C$};
  \fill (M) circle (1.5pt) node[below, font=\small] {$M$};
  \fill (N) circle (1.5pt) node[above left, font=\small] {$N$};
\end{tikzpicture}\hspace{1.1cm}%
\begin{tikzpicture}[scale=1.05]
  \coordinate (A) at (0,0);
  \coordinate (B) at (2.6,-1.3);
  \coordinate (C) at (1.4,-2.2);
  \coordinate (M) at ($(A)!-0.62!(B)$);
  \coordinate (N) at ($(A)!-0.62!(C)$);
  \draw[thick] (M) -- (B);
  \draw[thick] (N) -- (C);
  \draw[omThm, very thick] (M) -- (N);
  \draw[omThm, very thick] (B) -- (C);
  \fill (A) circle (1.5pt) node[right=2pt, font=\small] {$A$};
  \fill (B) circle (1.5pt) node[right, font=\small] {$B$};
  \fill (C) circle (1.5pt) node[below, font=\small] {$C$};
  \fill (M) circle (1.5pt) node[left, font=\small] {$M$};
  \fill (N) circle (1.5pt) node[above, font=\small] {$N$};
\end{tikzpicture}

{\small 泰勒斯定理的两种构型：嵌套三角形（左）与
“蝴蝶形”，其中 $M$ 与 $N$ 位于 $A$ 的另一侧（右）。
两图中均有 $(MN) \parallel (BC)$，且三个比值相等。}
\end{omfigure}

\begin{method}[用泰勒斯定理计算长度]\label{met:g9:thales:compute}
\begin{enumerate}
  \item 找出两直线相交的点 $A$ 以及两条
        \omterm{def:g6:lines:perp}{平行}线，并为两个三角形命名；
  \item 写出三个相等的比，且始终为“小三角形
        比大三角形”；
  \item 保留涉及三个已知长度与未知长度的两个比
        相等的关系；
  \item 通过交叉相乘解所得的比例式。
\end{enumerate}
\end{method}

\begin{example}\label{ex:g9:thales:compute}
在上图左侧的构型中，设 $AM = 3$，$AB = 5$，$AN = 2.4$
且 $MN = 3.3$，求 $AC$ 与 $BC$。由泰勒斯定理得
\[
\frac{AM}{AB} = \frac{AN}{AC} = \frac{MN}{BC},
\qquad\text{即}\qquad
\frac35 = \frac{2.4}{AC} = \frac{3.3}{BC}.
\]
由 $\frac35 = \frac{2.4}{AC}$，交叉相乘得
$3 \times AC = 5 \times 2.4 = 12$，所以 $AC = 4$。另由
$\frac35 = \frac{3.3}{BC}$ 得 $3 \times BC = 16.5$，所以 $BC = 5.5$。
\end{example}

\begin{example}[生活中的泰勒斯定理]\label{ex:g9:thales:stick}
一根高 $1$ m 的竖直木棍投下长 $1.5$ m 的影子，同一时刻一棵树
投下长 $12$ m 的影子。太阳光线\omterm{def:g6:lines:perp}{平行}，因此
木棍与影子所成的三角形、树与影子所成的三角形
处于泰勒斯构型：
\[
\frac{\text{树高}}{1} = \frac{12}{1.5},
\qquad\text{所以树高为 } 8 \text{ m。}
\]
\end{example}

\section{逆定理}

\begin{theorem}[泰勒斯定理的逆定理]\label{thm:g9:thales:converse}
点的放置方式与\cref{thm:g9:thales:direct}所示相同（$M$ 在 $(AB)$ 上，
$N$ 在 $(AC)$ 上，两条直线上的点顺序相同）：
若
\[
\frac{AM}{AB} = \frac{AN}{AC},
\]
则直线 $(MN)$ 与 $(BC)$ \omterm{def:g6:lines:perp}{平行}。
\end{theorem}
\admitted

\begin{method}[证明或否定平行关系]\label{met:g9:thales:converse}
\begin{enumerate}
  \item 分别计算两个比 $\dfrac{AM}{AB}$ 与
        $\dfrac{AN}{AC}$（保留为\omterm{def:g9:fractions:fraction}{分数}，不作\omterm{def:g6:decimals:rounding}{四舍五入}）；
  \item 若二者相等，\emph{且}两点在两条直线上的顺序相同，
        则两直线\omterm{def:g6:lines:perp}{平行}（泰勒斯定理的逆定理）；
  \item 若二者不相等，则两直线\emph{不}\omterm{def:g6:lines:perp}{平行}——因为若它们\omterm{def:g6:lines:perp}{平行}，
        由泰勒斯定理必有两个比相等。
\end{enumerate}
\end{method}

\begin{example}\label{ex:g9:thales:converse}
在第一条直线上，$AM = 4$，$AB = 10$；在第二条直线上，$AN = 6$，$AC = 15$。
于是
\[
\frac{AM}{AB} = \frac{4}{10} = \frac25,
\qquad
\frac{AN}{AC} = \frac{6}{15} = \frac25 :
\]
两比相等，且顺序相同，故 $(MN) \parallel (BC)$。

若改取 $AC = 14$，则 $\frac{6}{14} = \frac37 \neq \frac25$，
两直线不\omterm{def:g6:lines:perp}{平行}。
\end{example}

\section{放大与缩小}

\begin{definition}[图形的放缩]\label{def:g9:thales:scaling}
将图形按\emph{比例因子}\index{比例因子}
$k > 0$ 放大或缩小，指将其\emph{所有}长度乘以 $k$：
当 $k > 1$ 时为\emph{放大}，当
$k < 1$ 时为\emph{缩小}。在泰勒斯构型中，
三角形 $AMN$ 是三角形 $ABC$ 按比例因子 $k = \frac{AM}{AB}$ 的缩小图形。
\end{definition}

\begin{proposition}[对角度与面积的影响]\label{prop:g9:thales:scaling}
按\omterm{def:g9:thales:scaling}{比例因子} $k$ 放缩，角度与图形的形状保持不变，
每一长度乘以 $k$，并将每一\emph{\omterm{def:g6:measure:area}{面积}}乘以
$k^2$。
\end{proposition}
\admitted

\begin{omfigure}
\begin{tikzpicture}[scale=0.72]
  \draw[thick, omDef, fill=omDef!8] (0,0) rectangle (2,1.2);
  \draw[thick, omThm, fill=omThm!8] (3.4,0) rectangle (7.4,2.4);
  \node[font=\small, omDef] at (1,-0.45) {$2 \times 1.2$};
  \node[font=\small, omThm] at (5.4,-0.45) {$4 \times 2.4$};
  \node[font=\small] at (1,0.6) {面积 $2.4$};
  \node[font=\small] at (5.4,1.2) {面积 $9.6$};
\end{tikzpicture}

{\small 长度加倍（$k = 2$）时，\omterm{def:g6:measure:area}{面积}乘以
$k^2 = 4$：四个小长方形恰好拼成大长方形。}
\end{omfigure}

\begin{example}
一张 $10 \times 15$ cm 的照片按\omterm{def:g9:thales:scaling}{比例因子} $k = 3$ 放大，
冲洗后的尺寸为 $30 \times 45$ cm。其\omterm{def:g6:measure:area}{面积}由 $150$ cm$^2$ 变为
$150 \times 9 = 1350$ cm$^2$——是原来的 $9$ 倍，而非 $3$ 倍。
\end{example}

\section{练习}

\begin{exercise}[$\star$]\label{exo:g9:thales:1}
已知直线 $(MN)$ 与 $(BC)$ \omterm{def:g6:lines:perp}{平行}，$M$ 在 $[AB]$ 上，$N$ 在
$[AC]$ 上；$AM = 2$，$AB = 6$，$AN = 3$，$BC = 9$。求 $AC$ 与 $MN$。
\end{exercise}

\begin{exercise}[$\star$]\label{exo:g9:thales:2}
同一构型：$AM = 5$，$MB = 3$（注意是 $MB$，而非 $AB$！），
$AN = 4$。求 $AC$。
\end{exercise}

\begin{exercise}[$\star$]\label{exo:g9:thales:3}
在蝴蝶型构型中，$A$ 介于 $M$ 与 $B$ 之间，又介于 $N$
与 $C$ 之间，且 $(MN) \parallel (BC)$；$AM = 3$，$AB = 7.5$，$AN = 2$，
$MN = 2.6$。求 $AC$ 与 $BC$。
\end{exercise}

\begin{exercise}[$\star$]\label{exo:g9:thales:4}
$M$ 在 $[AB]$ 上，$AM = 6$，$AB = 8$；$N$ 在 $[AC]$ 上，$AN = 9$，
$AC = 12$。直线 $(MN)$ 与 $(BC)$ 是否\omterm{def:g6:lines:perp}{平行}？
\end{exercise}

\begin{exercise}[$\star\star$]\label{exo:g9:thales:5}
与上题条件相同，但 $AM = 4$，$AB = 6$，$AN = 5$，$AC = 8$。直线 $(MN)$
与 $(BC)$ 是否\omterm{def:g6:lines:perp}{平行}？请仔细说明理由。
\end{exercise}

\begin{exercise}[$\star\star$]\label{exo:g9:thales:6}
一人身高 $1.8$ m，站在距路灯底部 $2$ m 处；
其影长为 $3$ m。画出由路灯、此人及影子顶端
构成的泰勒斯构型，
并求路灯的高度。
\end{exercise}

\begin{exercise}[$\star\star$]\label{exo:g9:thales:7}
某地图的\omterm{def:g7:prop:scale}{比例尺}为 $1 : 25\,000$（地图上的长度等于实际长度
乘以 $k = \frac{1}{25000}$）。
\begin{enumerate}
  \item 地图上两村庄相距 $6.8$ cm，其实际距离
        是多少 km？
  \item 某森林在地图上的\omterm{def:g6:measure:area}{面积}为 $8$ cm$^2$。其实际\omterm{def:g6:measure:area}{面积}
        是多少 km$^2$？
\end{enumerate}
\end{exercise}

\begin{exercise}[$\star\star$]\label{exo:g9:thales:8}
三角形 $ABC$ 中，$AB = 12$，$AC = 15$，$BC = 18$。点 $M$ 在
$[AB]$ 上，满足 $AM = 8$；过 $M$ 且\omterm{def:g6:lines:perp}{平行}于 $(BC)$ 的直线
与 $[AC]$ 相交于 $N$。
\begin{enumerate}
  \item 求 $AN$ 与 $MN$。
  \item 由 $ABC$ 到 $AMN$ 的\omterm{def:g9:thales:scaling}{比例因子}是多少？比较
        两三角形的\omterm{def:g6:measure:perimeter}{周长}。
\end{enumerate}
\end{exercise}

\begin{exercise}[$\star\star\star$]\label{exo:g9:thales:9}
设 $ABCD$ 为梯形，$(AB) \parallel (CD)$，$AB = 4$，且
$CD = 6$，其对角线 $[AC]$ 与 $[BD]$ 相交于 $O$。
利用以 $O$ 为中心的蝴蝶型构型中的泰勒斯定理，求
比 $\dfrac{OA}{OC}$，并证明 $O$ 将两条对角线分成
相同的比。
\end{exercise}

\section{问题：无刻度尺、针孔相机与
梯形的三种平均数}

\begin{problem}[{周末问题——泰勒斯定理的应用：任意
线段等分、用鞋盒测量太阳，以及
三种著名平均数汇聚的梯形}]
\label{pb:g9:thales:1}
泰勒斯定理研究的是\emph{射线}：阳光的射线、
透过针孔的光线、从\omterm{def:g5:solids:def}{顶点}发出的射线束。
本题从三个方面运用它：作为作图工具（将
\emph{任意}长度的\omterm{def:g6:lines:objects}{线段}（甚至 $\sqrt2$）分成完全相等的若干段）、
作为测量仪器（用鞋盒和硬币测量太阳的\omterm{def:g4:geometry:circle}{直径}），
以及作为放大镜，观察
\cref{exo:g9:thales:9}中的梯形，其中
本系列所讲的算术、几何与\omterm{pb:g8:speed:1}{调和平均数}全部出现在同一幅图中。

\medskip
\noindent\textbf{第一部分——用无刻度尺等分\omterm{def:g6:lines:objects}{线段}。}
\begin{enumerate}
  \item 画一条\omterm{def:g6:lines:objects}{线段} $[AB]$，长 $7$ cm。
        用圆规和无刻度尺将其三等分：先任作一条射线，
        过 $A$ 作射线（不经过 $B$），在其上截取三段
        相等的圆规长度，得点 $P_1$、$P_2$、$P_3$；
        连接 $P_3$ 与 $B$，作 $(P_3 B)$ 的\omterm{def:g6:lines:perp}{平行}线，
        分别过 $P_1$、$P_2$。完成作图，标出
        \omterm{def:g6:lines:perp}{平行}线与 $[AB]$ 相交的位置。
  \item 由\cref{thm:g9:thales:direct}可知，
        所标出的两点恰好将 $[AB]$ 三等分。
  \item 改用此法，将\omterm{def:g6:lines:objects}{线段}分成 $5$ 等份，
        再作出已知\omterm{def:g6:lines:objects}{线段}的 $\frac35$。
  \item 作点 $M$，使其位于 $[AB]$ 上，且
        $\frac{AM}{MB} = \frac23$（即 $M$ 位于全长的 $\frac25$ 处，
        从 $A$ 量起）。射线上应截取几段相等长度？
        应把哪一点与 $B$ 相连？
  \item 为什么这种作图法比用有刻度的尺测量更好？
        考虑长为 $\sqrt2$ 的\omterm{def:g6:lines:objects}{线段}，
        它是单位正方形的对角线，由
        \cref{pb:g9:sqrt:1}可知它是无理数。若用尺测量会得到什么，
        而用泰勒斯定理作图又得到什么？
\end{enumerate}

\medskip
\noindent\textbf{第二部分——针孔相机。}
在一个封闭箱子的一面上扎一个小孔，在对面
出现外界景物的倒立像。物体上的一点、小孔与像点
三点共线——光沿直线传播——因此
物体（高为 $H$，距小孔前方 $D$ 处）
与其像（高为 $h$，位于小孔后方
深度为 $d$ 的后墙上）
构成泰勒斯定理的蝴蝶型构型。
\begin{enumerate}[resume]
  \item 由泰勒斯定理推出针孔公式
        $h = H \times \frac dD$，
        其中用到小孔周围的蝴蝶型构型。
  \item 一棵 $12$ m 高的树距鞋盒 $30$ m，
        鞋盒深 $20$ cm。它的像有多高？是正立还是倒立？
  \item 将鞋盒对准太阳：在小孔后方 $d = 1$ m 处，太阳的像
        是一个\omterm{def:g4:geometry:circle}{直径}约为
        $9$ mm 的圆盘。已知太阳距离为
        $D = 1.5 \times 10^{11}$ m，求其\omterm{def:g4:geometry:circle}{直径}，用\omterm{def:g9:fractions:scientific}{科学记数法}表示
        （运用\cref{pb:g9:fractions:1}中的
        记号）。将结果与真实值作比较，
        真实值为 $1.39 \times 10^9$ m。
  \item 月球的\omterm{def:g4:geometry:circle}{直径}为 $3.5 \times 10^6$ m，距离为
        $3.8 \times 10^8$ m。求比值
        $\frac{\text{直径}}{\text{距离}}$，分别求月球和
        太阳的这一比值。这两个数揭示了怎样的巧合——
        它又使哪种壮观的天象
        成为可能？
  \item 用一两句话说明：为什么针孔像是倒立的，
        以及为什么扩大针孔会使像变亮，
        却也更加模糊？
\end{enumerate}

\medskip
\noindent\textbf{第三部分——三种平均数的
梯形。}
$ABCD$ 即\cref{exo:g9:thales:9}中的梯形：
$(AB) \parallel (CD)$，$AB = 4$，$CD = 6$，对角线相交于
$O$，且 $\frac{OA}{OC} = \frac{OB}{OD} = \frac{4}{6} =
\frac23$。
\begin{enumerate}[resume]
  \item 蝴蝶型练习：两条直线相交于 $O$，
        $(MN) \parallel (PQ)$ 构成蝴蝶型，$OM = 3$，
        $OP = 7.5$，$MN = 4$，$OQ = 6$。求 $PQ$ 与
        $ON$。
  \item 说明：以比 $\frac12$ 应用泰勒斯定理时，
        它如何包含
        \cref{thm:g8:midpoints:direct}作为其特殊情形。
  \item 过 $O$ 作与两底\omterm{def:g6:lines:perp}{平行}的直线，它与
        $[AD]$ 相交于 $E$。在三角形 $ACD$ 中（注意
        $\frac{AO}{AC} = \frac25$），求 $EO$。
  \item 仿照三角形 $BDC$ 中的对称计算，
        \omterm{def:g6:lines:objects}{线段} $OF$（$F$ 在 $[BC]$ 上）的长度与之相同。
        由此可得
        \[
        EF = 2 \times \frac{4 \times 6}{4 + 6} = 4.8 :
        \]
        过对角线交点的\omterm{def:g6:lines:perp}{平行}线量出了两底的
        \emph{\omterm{pb:g8:speed:1}{调和平均数}}——
        即\cref{pb:g8:speed:1}中的往返平均数，
        在纯几何中再次出现。
  \item 压轴题：求两底的三种经典平均数，
        即 $4$ 与 $6$ 的算术平均数
        $\frac{4+6}{2}$、\omterm{pb:g8:cosine:1}{几何平均数} $\sqrt{4 \times 6}$、\omterm{pb:g8:speed:1}{调和平均数}
        $\frac{2 \times 4 \times 6}{4 + 6}$，并将它们排序。
        三者都出现在梯形之中：算术平均数是
        连接两腰中点的中位线
        （\cref{pb:g7:areas:1}），\omterm{pb:g8:speed:1}{调和平均数}则是
        你所求得的\omterm{def:g6:lines:objects}{线段} $EF$，\omterm{pb:g8:cosine:1}{几何平均数}则是
        将梯形分成两个\emph{相似}梯形的\omterm{def:g6:lines:perp}{平行}线
        （可用计算器检验其数值；
        其证明是一道有趣的附加挑战）。这三种平均数之间的
        不等式，本书在哪些地方给出过证明
        （\cref{pb:g8:cosine:1}、\cref{pb:g8:speed:1}）？
\end{enumerate}
\end{problem}
